% Generated draft: requires mathematical and source review.
\documentclass[11pt]{article}
\usepackage[margin=25mm]{geometry}
\usepackage{fontspec}
\setmainfont{Latin Modern Roman}
\usepackage{amsmath,amssymb,mathtools,unicode-math}
\setmathfont{Latin Modern Math}
\usepackage{enumitem,needspace,xurl,hyperref}
\setlength{\parindent}{0pt}
\setlength{\parskip}{5pt}
\setlength{\emergencystretch}{3em}
\newcommand{\N}{\mathbb{N}}
\newcommand{\Z}{\mathbb{Z}}
\newcommand{\Q}{\mathbb{Q}}
\newcommand{\R}{\mathbb{R}}
\newcommand{\C}{\mathbb{C}}
\newcommand{\F}{\mathbb{F}}
\newcommand{\A}{\mathbb{A}}
\newcommand{\PP}{\mathbb{P}}
\newcommand{\E}{\mathbb{E}}
\newcommand{\eps}{\varepsilon}
\newcommand{\dd}{\,\mathrm{d}}
\begin{document}
\section*{UnsolvedMath ID: 20001945}
\textbf{Source problem code:} catalogue entry
\subsection{Definitions and mathematical statement}
\textbf{Question 4.} Identify the purported limits of the Cohn--Miller sequences, explain the observed rational Taylor coefficients, and determine whether rational coefficients occur beyond the quadratic terms. The extract does not define the sequences, dimensions, normalization, root modification, or convergence topology. Consequently no unique concrete sequence or numerical coefficient list can be recovered from this record alone.

\textbf{Conditional formulation of the supplied aid.} For $d\in\{8,24\}$, write $f_{d,m}$ for the intended normalized modified-root Cohn--Miller functions and $M_d$ for the corresponding normalized magic function, leaving their omitted construction explicit. The radial Schwartz space $\mathcal S_{\mathrm{rad}}(\R^d)$ consists of smooth radial functions with all seminorms
\[
 p_{A,\alpha}(f)=\sup_{x\in\R^d}(1+|x|)^A|\partial^\alpha f(x)|
\]
finite, for every nonnegative integer $A$ and multi-index $\alpha$. A cluster point means a limit of a subsequence with indices tending to infinity in this topology. The aid asserts
\[
 f_{d,m_j}\longrightarrow f\text{ in }\mathcal S_{\mathrm{rad}}
 \quad\Longrightarrow\quad f=M_d.
\]
If the entire sequence is relatively compact in this space, this uniqueness of cluster points implies $f_{d,m}\to M_d$. The construction of $M_d$ independently of the sequence does not establish relative compactness or convergence of the sequence. Nor is a unique cluster-point statement an existence statement for a cluster point.

To make the aid's Taylor--moment language precise under an explicit convention, put
\[
 \widehat f(\xi)=\int_{\R^d}f(x)e^{-2\pi i x\cdot\xi}\,\dd x.
\]
For radial Schwartz $f$, the coefficient of the finite Taylor jet of $r\mapsto\widehat f(re_1)$ of order $2j$ is
\[
 \frac1{(2j)!}\left.\frac{\dd^{2j}}{\dd r^{2j}}
        \widehat f(re_1)\right|_{r=0}
 =\frac{(-1)^j\pi^{2j}}{j!(d/2)_j}
       \int_{\R^d}|x|^{2j}f(x)\,\dd x,
\]
where $(a)_0=1$ and $(a)_j=a(a+1)\cdots(a+j-1)$. In particular the quadratic coefficient is $-2\pi^2/d$ times the second moment. The analogous identity for $f$ uses $\widehat f$, since radial functions are even. This is an explicitly chosen Fourier normalization, not a recovery of an omitted source convention. Schwartz convergence transfers each finite Taylor jet and each polynomially weighted moment; it need not imply convergence of an infinite Taylor series to a Schwartz function. The aid reports that all four conjectured quadratic coefficients for the dimension-eight and dimension-twenty-four functions and their transforms are proved, but supplies neither their values nor their proving citations here. It reports no verified theorem for additional rational coefficients or for irrationality of higher coefficients. No such values or theorem are invented.

\textbf{Adjacent Epstein-zeta context.} For a full-rank covolume-one lattice $L\subset\R^n$, put
\[
 \theta_L(t)=\sum_{v\in L}e^{-\pi t|v|^2},\qquad
 E(L,s)=\sum_{v\in L\setminus\{0\}}|v|^{-2s}
 \quad(\operatorname{Re}s>n/2).
\]
Then the usual Mellin relation is
\[
 \pi^{-s}\Gamma(s)E(L,s)
 =\int_0^\infty(\theta_L(t)-1)t^{s-1}\,\dd t
 \quad(\operatorname{Re}s>n/2).
\]
These objects depend on the vector-length spectrum with multiplicities. The source attributes continuation to Epstein; ``for all $s$'' must be understood meromorphically, with a pole at $s=n/2$, not as a finite value there. This paragraph supplies context, not a definition of the Cohn--Miller sequences or an additional zeta-minimization question. The extract contains the contributor line Peter Sarnak and Andreas Str\"ombergsson, whose precise attachment to the question or ensuing context should be checked.
\subsection{Short English statement}
The source asks what the Cohn--Miller approximations converge to and why some Taylor coefficients are rational. The aid identifies possible Schwartz limits in dimensions eight and twenty-four and reduces convergence to relative compactness. The actual sequence construction, normalizations, four quadratic values and proving references are missing; higher rationality remains unverified.
\subsection{Sources}
\begin{itemize}
\item Local UnsolvedMath registry, permanent ID 20001945, source code catalogue entry; canonical \texttt{aim-geometry-notes.json}, notes[282].
\item AIM workshop \emph{Sphere packings, lattices, and infinite dimensional algebra}, Question 4 and adjacent Epstein-zeta paragraph: \url{https://aimath.org/WWN/spherepacking/spherepacking.pdf}.
\item Contributor line in the supplied extract: Peter Sarnak and Andreas Str\"ombergsson. Analytic continuation of the lattice zeta function is attributed to Epstein. No separate authoritative citations defining the sequences, constructing the magic functions or proving the four quadratic values are supplied.
\item Registry research attempt 1, dated 2026-08-14, partial result, rigor ``proved'', novelty confidence low. The aid is explicitly machine-generated, not peer reviewed, and requires independent verification before citation.
\end{itemize}
\end{document}
